\section{Условие}
{\bfseries Цель работы:} 
Приобретение практических навыков в:
\begin{itemize}
  \item Управление процессами в ОС
  \item Обеспечение обмена данных между процессами посредством каналов
\end{itemize}

{\bfseries Задание:} Составить и отладить программу на языке Си, осуществляющую работу с процессами и
взаимодействие между ними в одной из двух операционных систем. В результате работы
программа (основной процесс) должен создать для решение задачи один или несколько
дочерних процессов. Взаимодействие между процессами осуществляется через системные
сигналы/события и/или каналы (pipe).
Необходимо обрабатывать системные ошибки, которые могут возникнуть в результате работы.

{\bfseries Вариант:} 14

{\bfseries Задание варианта:}
Родительский процесс создает два дочерних процесса. Перенаправление стандартных потоков
ввода-вывода показано на картинке выше. Child1 и Child2 можно «соединить» между собой
дополнительным каналом. Родительский и дочерний процесс должны быть представлены
разными программами.
Родительский процесс принимает от пользователя строки произвольной длины и пересылает их в
pipe1. Процесс child1 и child2 производят работу над строками. Child2 пересылает результат своей
работы родительскому процессу. Родительский процесс полученный результат выводит в
стандартный поток вывода.
Child1 переводит строки в нижний регистр. Child2 убирает все задвоенные пробелы

